\documentclass[11pt,twoside]{article}

\addtolength{\textwidth}{0.5in}
\usepackage{epsfig,amsfonts,color}
%\usepackage[leqno]{amsmath}
\usepackage{amsmath}
\usepackage{optidef}
%\usepackage{showlabels}
\newcommand{\norm}[1]{\left| \left| #1 \right| \right|}
\newcommand{\mods}[1]{\left|  #1 \right|}
\newcommand{\cm}[1]{{\color{black} #1}}
\newcommand{\cv}[1]{{\color{black} #1}}
\usepackage{amssymb, palatino, geometry,url}
\usepackage{algorithmic}
\usepackage{subcaption}
\usepackage{graphicx}
\graphicspath{ {./Figures/} }
\usepackage{longtable}
\usepackage[section]{placeins}
\usepackage[noresetcount,lined,boxed]{algorithm2e} % ... for algorithms
\newcommand{\mynote}[1]{{ \color{black}{#1}}}
\usepackage[colorlinks=true,linkcolor=blue,citecolor=blue,urlcolor=blue]{hyperref}
%\usepackage{cite}

\usepackage{indentfirst}
\usepackage{parskip}
\setlength{\parindent}{12pt}

\makeatletter
\newcommand{\nocaption}{\addtocounter{\@captype}{-1}} 
% stop counting the figure numbering when having no caption command
\makeatother
\pagenumbering{arabic}

\geometry{letterpaper,
	left       = 0.9in,
	right      = 0.9in,
	top        = 0.9in,
	bottom     = 0.9in}
\linespread{1.2}

\usepackage{fancyhdr}
\pagestyle{fancy}


\newtheorem{assumption}{Assumption}
% ... number equations within sections
\renewcommand{\theequation}{\thesection.\arabic{equation}}

% ... mathematical symbols / macros
\newcommand{\R}{\mbox{I}\!\mbox{R}}
\newcommand{\Tau}{\mathcal{T}}
\newcommand{\C}{\mbox{I}\hspace{-0.5em}\mbox{C}}
%\newcommand{\mini}{\mathop{\mbox{minimize}}}
%\newcommand{\maxi}{\mathop{\mbox{maximize}}}
\newcommand{\amin}{\mathop{\mbox{argmin}}}
\newcommand{\amax}{\mathop{\mbox{argmax}}}
\newcommand{\st}{\mbox{subject to }}
\newcommand{\bmath}[1]{\mbox{\boldmath $ #1 $}}
\newcommand{\be}{\begin{equation}}
	\newcommand{\ee}{\end{equation}}
\newcommand{\bea}{\begin{eqnarray}}
	\newcommand{\eea}{\end{eqnarray}}
\newcommand{\dps}{\displaystyle}
\newcommand{\mbs}[1]{\mbox{\scriptsize #1}}
\newcommand{\pref}[1]{(\ref{#1})}
\newcommand{\bvec}{\left(\begin{array}{c}}
	\newcommand{\evec}{\end{array}\right)}
\newcommand{\bsub}{\begin{subequations}}
	\newcommand{\esub}{\end{subequations}}
\newcommand{\Model}{{\cal M}}
\newcommand{\ModelTLM}{{\mathbf{M}}}
\newcommand{\ModelADJ}{{\mathbf{M}^T}}
\newcommand{\ModelCovMat}{\mathbf{Q}}
\newcommand{\CholeskyMatrix}{\mathbf{L}}

\newcommand{\ti}{{t_0}}
\newcommand{\tf}{{t_F}}
\newcommand{\tk}[1]{{t_{#1}}}
\newcommand{\xnarr}{{x_{\scriptscriptstyle{\rm{NARR}}}}}

\newcommand{\CostFcn}{{\Psi}}
\newcommand{\SampleVarMat}{{\mathbf{S}^2}}
\newcommand{\VarMat}{\mathbf{D}}
\newcommand{\CorMat}{\mathbf{C}}
\newcommand{\CovMat}{\mathbf{V}}
\newcommand{\Avg}{\mu}
\newcommand{\NEnsembleWindows}{N}
\newcommand{\NEnsembleMembers}{{N_{S}}}
\newcommand{\NState}{M_x}
\newcommand{\NSteps}{m}
\newcommand{\Expect}[1]{\overline{#1}}
\newcommand{\mean}{{\rm mean}}
\newcommand{\Lagr}{\mathcal{L}}
\usepackage{lineno}
\linenumbers


\begin{document}
	
\title{Decision-Focused NLP Surrogates for Convex MINLPs}

\author{Parth Brahmbhatt, Kiernan Jennings\\
	{\small Department of Chemical and Biological Engineering}\\
	{\small \;University of Wisconsin, 1415 Engineering Dr, Madison, WI 53706, USA}}
\date{}
\maketitle
	
	
\begin{abstract}	
	Convex Mixed-Integer Nonlinear Programs (MINLPs) are crucial for modeling complex engineering and energy systems. However, solving them relies on computationally exhaustive branch-and-bound algorithms, making them largely unsuitable for real-time, online applications. Recent research has explored decision-focused surrogate modeling to accelerate combinatorial optimization, typically by formulating Linear Programming (LP) surrogates for Mixed-Integer Linear Programs (MILPs) [1] or large LPs [2] using bilevel programming. Concurrently, there have been significant advances in integrating differentiable optimization layers directly into machine learning pipelines [3]. However, a major gap remains in extending these techniques to nonlinear domains. Specifically, there is a distinct need to develop continuous Nonlinear Programming (NLP) surrogates for complex Mixed-Integer Nonlinear Programs (MINLPs) by learning the feasible region constraints in a decision-focused manner.
	
	This project will develop an end-to-end differentiable NLP surrogate to rapidly solve convex MINLPs. We will train a neural network to learn problem-specific constraints that guide the continuous relaxation, followed by a feasibility restoration step to quickly recover valid discrete decisions. To demonstrate the effectiveness and speed of this architecture, we will showcase a case study focusing on a hybrid vehicle control problem or a Unit Commitment AC Optimal Power Flow (UC-ACOPF) problem, formulated as a convex Mixed-Integer Quadratically Constrained Quadratic Program (MIQCQP) relaxation [4].
	
%	[1] Dixit, S., Gupta, R., & Zhang, Q. (2024). Decision-focused surrogate modeling for mixed-integer linear optimization. arXiv preprint arXiv:2406.05697.
%	
%	[2] Anrrango, A., Quisaguano, A., Constante-Flores, G. E., & Li, C. (2026). Self-Supervised Learning of Parametric Approximation for Security-Constrained DC-OPF. arXiv preprint arXiv:2601.13486.
%	
%	[3] Agrawal, A., Amos, B., Barratt, S., Boyd, S., Diamond, S., & Kolter, J. Z. (2019). Differentiable convex optimization layers. Advances in neural information processing systems, 32.
%	
%	[4] Constante-Flores, G., & Li, C. (2026). A quadratically-constrained convex approximation for the AC optimal power flow. Optimization and Engineering, 1-23.
	
\end{abstract}
	
\section{Concept}

We are trying to learn the following convex MIQCQP:
\begin{equation}
	\begin{aligned}
		P(u) = \quad & \min_{x} & & c(u)^\top x \\
		& \text{s.t.} & & x^\top Q(u) x + R(u)^\top x \leq 0 \\
		&&& x \in \mathbb{R}^m \times \mathbb{Z}^{n-m}	
	\end{aligned}
\end{equation}


Using the following surrogate convex QCQP:
\begin{equation}
	\begin{aligned}
		\hat{P}(u,\theta) = \quad & \min_{x} & & c(u)^\top x \\
		& \text{s.t.} & & x^\top Q(u) x + R(u)^\top x \leq 0 \\
		&&& \tilde{A}(u, \theta) x \leq \tilde{b}(u, \theta) \\
		&&& x \in \mathbb{R}^n
	\end{aligned}
\end{equation}

We are learning convex linearization points and linear cuts for the convex MIQCQP (UC-ACOPF) to produce a surrogate convex QCQP. In cases of infeasible integer solution points, we project the solution using $\ell_1$-norm minimization for feasibility restoration.

\section{Preliminaries}
Theory:
\begin{equation}
	\begin{aligned}
		QP = \quad & \min_{x\in\mathbb{R}^n} & &\frac{1}{2} x^\top Q x + c^\top x \\
		& \text{s.t.} & & Gx \leq h & (\lambda) \\
		&&& Ax = b & (\mu) 
	\end{aligned}
\end{equation}

QP KKT:
\begin{equation}
	\begin{aligned}
		Qx + c +A^\top\mu + G^\top\lambda = 0 & & (\nabla \Lagr)\\
		Ax=b && (P_{eq}) \\
		\lambda_i (h-Gx)_i = 0 & \quad \forall i & (C) 
	\end{aligned}
\end{equation}
Differentiate and translate to matrix form:
\begin{equation}
	\begin{bmatrix}
		Q & G^\top & A^\top \\
		D(\lambda)G & D(h - Gx) & 0 \\
		A & 0 & 0
	\end{bmatrix}
	\begin{bmatrix}
		dx \\
		d\lambda \\
		d\mu
	\end{bmatrix}
	= -
	\begin{bmatrix}
		dQ \, x + dc + dG^\top \lambda + dA^\top \mu \\
		D(\lambda)dh - D(\lambda)dG \, x \\
		dAx - db
	\end{bmatrix},
\end{equation}
where $D(\cdot)$ is the diagonal matrix of a given vector.

\section{Unit Commitment with Alternating Current Optimal Power Flow}
The lifted rectangular coordinates formulation for the AC OPF problem,
\begin{subequations}
	\begin{align}
		& \min_{\Xi} \sum_{k \in \mathcal{G}} C_k(p_k^{\text{g}}) \tag{1a} \\
		& \text{subject to} \notag \\
		& \sum_{k \in \mathcal{G}_i} p_k^{\text{g}} - P_i^{\text{d}} - G_i^{\text{sh}} c_{ii} = \sum_{(i,j) \in \mathcal{E}_i} p_{ij} + \sum_{(j,i) \in \mathcal{E}_i} p_{ji}, \forall i \in \mathcal{N}, \tag{1b} \\
		& \sum_{k \in \mathcal{G}_i} q_k^{\text{g}} - Q_i^{\text{d}} + B_i^{\text{sh}} c_{ii} = \sum_{(i,j) \in \mathcal{E}_i} q_{ij} + \sum_{(j,i) \in \mathcal{E}_i} q_{ji}, \forall i \in \mathcal{N}, \tag{1c} \\
		& p_{ij} = \frac{G_{ij}}{T_{ij}^2} c_{ii} - \frac{G_{ij}}{T_{ij}} c_{ij} + \frac{B_{ij}}{T_{ij}} s_{ij}, \forall (i, j) \in \mathcal{E}, \tag{1d} \\
		& p_{ji} = G_{ij} c_{jj} - \frac{G_{ij}}{T_{ij}} c_{ij} - \frac{B_{ij}}{T_{ij}} s_{ij}, \forall (i, j) \in \mathcal{E}, \tag{1e} \\
		& q_{ij} = -\frac{B_{ij} + B_{ij}^{\text{sh}}}{T_{ij}^2} c_{ii} + \frac{G_{ij}}{T_{ij}} s_{ij} + \frac{B_{ij}}{T_{ij}} c_{ij}, \forall (i, j) \in \mathcal{E}, \tag{1f} \\
		& q_{ji} = -\left( B_{ij} + B_{ij}^{\text{sh}} \right) c_{jj} - \frac{G_{ij}}{T_{ij}} s_{ij} + \frac{B_{ij}}{T_{ij}} c_{ij}, \forall (i, j) \in \mathcal{E}, \tag{1g} \\
		& (p_{ij})^2 + (q_{ij})^2 \leq \overline{S}_{ij}^2, \forall (i, j), (j, i) \in \mathcal{E}, \tag{1h} \\
		& \underline{V}_i^2 \leq c_{ii} \leq \overline{V}_i^2, \forall i \in \mathcal{N}, \tag{1i} \\
		& \underline{P}_k \leq p_k^{\text{g}} \leq \overline{P}_k, \forall k \in \mathcal{G}, \tag{1j} \\
		& \underline{Q}_k \leq q_k^{\text{g}} \leq \overline{Q}_k, \forall k \in \mathcal{G}, \tag{1k} \\
		& c_{ii} = (v_i^{\text{re}})^2 + (v_i^{\text{im}})^2, \forall i \in \mathcal{N}, \tag{1l} \\
		& c_{ij} = v_i^{\text{re}} v_j^{\text{re}} + v_i^{\text{im}} v_j^{\text{im}}, \forall (i, j) \in \mathcal{E}, \tag{1m} \\
		& s_{ij} = v_i^{\text{re}} v_j^{\text{im}} - v_j^{\text{re}} v_i^{\text{im}}, \forall (i, j) \in \mathcal{E}, \tag{1n}
	\end{align}
\end{subequations}

Using Constante-Flores and Li 2026 convex formulation,
\begin{subequations}
	\begin{align}
		& \min_{\Xi} \sum_{k \in \mathcal{G}} C_k(p_k^{\text{g}}) + \rho \left[ \sum_{i \in \mathcal{N}} \xi_i^c + \sum_{(i,j) \in \mathcal{E}} (\xi_{ij}^c + \xi_{ij}^s) \right] \\
		& \text{subject to} \quad (1\text{b}) - (1\text{k}), \notag \\
		& c_{ii} \geq (v_i^{\text{re}})^2 + (v_i^{\text{im}})^2, \forall i \in \mathcal{N} \\
		& c_{ii} \leq 2 \left( V_i^{\text{re}} v_i^{\text{re}} + V_i^{\text{im}} v_i^{\text{im}} \right) - \left( (V_i^{\text{re}})^2 + (V_i^{\text{im}})^2 \right) + \xi_i^c, \forall i \in \mathcal{N}\\
		& (v_i^{\text{re}} + v_j^{\text{re}})^2 + (v_i^{\text{im}} + v_j^{\text{im}})^2 - 4c_{ij} \leq \xi_{ij}^c + \notag \\
		& \quad 2(v_i^{\text{re}} - v_j^{\text{re}})(V_i^{\text{re}} - V_j^{\text{re}}) + 2(v_i^{\text{im}} - v_j^{\text{im}})(V_i^{\text{im}} - V_j^{\text{im}}) - \notag \\
		& \quad \left[ (V_i^{\text{re}} - V_j^{\text{re}})^2 + (V_i^{\text{im}} - V_j^{\text{im}})^2 \right], \forall (i, j) \in \mathcal{E}, \\
		& (v_i^{\text{re}} - v_j^{\text{re}})^2 + (v_i^{\text{im}} - v_j^{\text{im}})^2 + 4c_{ij} \leq \xi_{ij}^c + \notag \\
		& \quad 2(v_i^{\text{re}} + v_j^{\text{re}})(V_i^{\text{re}} + V_j^{\text{re}}) + 2(v_i^{\text{im}} + v_j^{\text{im}})(V_i^{\text{im}} + V_j^{\text{im}}) - \notag \\
		& \quad \left[ (V_i^{\text{re}} + V_j^{\text{re}})^2 + (V_i^{\text{im}} + V_j^{\text{im}})^2 \right], \forall (i, j) \in \mathcal{E} \\
		& (v_i^{\text{re}} - v_j^{\text{im}})^2 + (v_j^{\text{re}} + v_i^{\text{im}})^2 + 4s_{ij} \leq \xi_{ij}^s + \notag \\
		& \quad 2(v_i^{\text{re}} + v_j^{\text{im}})(V_i^{\text{re}} + V_j^{\text{im}}) + 2(v_j^{\text{re}} - v_i^{\text{im}})(V_j^{\text{re}} - V_i^{\text{im}}) - \notag \\
		& \quad \left[ (V_i^{\text{re}} + V_j^{\text{im}})^2 + (V_j^{\text{re}} - v_i^{\text{im}})^2 \right], \forall (i, j) \in \mathcal{E},\\
		& (v_i^{\text{re}} + v_j^{\text{im}})^2 + (v_j^{\text{re}} - v_i^{\text{im}})^2 - 4s_{ij} \leq \xi_{ij}^s + \notag \\
		& \quad 2(v_i^{\text{re}} - v_j^{\text{im}})(V_i^{\text{re}} - V_j^{\text{im}}) + 2(v_j^{\text{re}} + v_i^{\text{im}})(V_j^{\text{re}} + V_i^{\text{im}}) - \notag \\
		& \quad \left[ (V_i^{\text{re}} - V_j^{\text{im}})^2 + (V_j^{\text{re}} + V_i^{\text{im}})^2 \right], \forall (i, j) \in \mathcal{E},  \\
		& \xi_i^c \geq 0, \forall i \in \mathcal{N},\\
		& \xi_{ij}^c \geq 0, \forall (i, j) \in \mathcal{E},  \\
		& \xi_{ij}^s \geq 0, \forall (i, j) \in \mathcal{E},
	\end{align}
\end{subequations}


\end{document}
